The elicitation of usability requirements remains a recurring challenge in Software Engineering,
particularly due to the difficulty of making usability aspects explicit during early development
phases. The USARP (Usability Requirements with Personas and User Stories) method was
proposed to support the systematic identification, discussion, and specification of usability
requirements through the integrated use of personas, user stories, and usability mechanisms
represented by cards. Although previous empirical studies have reported benefits related to the
method, an integrated synthesis of the available evidence across different contexts is still lacking.
This work presents an empirical synthesis of the USARP method, grounded in the integration of
quantitative and qualitative evidence from four primary studies conducted in academic and
industrial contexts. The methodology combines descriptive statistical analysis, dimensionality
reduction through Principal Component Analysis (PCA), and thematic analysis of open-ended
responses, consolidated at the respondent level, enabling the investigation of latent dimensions
associated with the perceived value of the method and the perceived usefulness of its mechanisms.
The results indicate a broadly positive and homogeneous perception of the method across academic
and industrial contexts, with a strong ceiling effect in the assessed items. The Principal Component
Analysis, of an exploratory nature, revealed a dimension of general positive perception and
secondary contrasts among the method's card blocks. Comparisons between contexts revealed no
statistically robust differences after correction for multiple comparisons, except for one item related
to the usability mechanism context. The qualitative analysis of the open-ended responses identified
the cards and specification guides as the main element supporting elicitation, along with tensions
related to the learning curve and the volume of cards, and the underuse of the prototyping cards. As
a contribution, this study provides a critical empirical synthesis of the USARP method, outlining its
strengths, limitations, and implications for its adoption in real-world software development settings.

\keywords{USARP; usability requirements; empirical synthesis; software engineering; personas and user stories.}